\chapter{System Analysis and Requirements}
\label{ch:requirements}

\section{Analysis of the Existing Workflow}

\noindent Before design, the current workflow in a typical college was
analysed and the following recurring weaknesses identified:

\begin{itemize}
  \item \textbf{Fragmented material.} Course material exists as scattered
  PDFs and scans that are hard to search, reuse, or convert into other
  formats.
  \item \textbf{No syllabus structure.} Syllabus documents are read manually
  and translated into teaching plans repeatedly for each batch.
  \item \textbf{Uncontrolled question creation.} Questions are composed ad hoc
  without tracking origin, approval, or alignment to the examination pattern.
  \item \textbf{Manual assessment.} Distribution, collection, and evaluation
  are manual, giving students delayed and summary-only feedback.
  \item \textbf{No access discipline.} Generic shared logins mean no
  accountability and no restriction of what a user may do or where.
\end{itemize}

\section{Functional Requirements}

\noindent The functional requirements of the platform are grouped by area. Each
requirement is implemented and verified (see Chapter~\ref{ch:testing}).

\begin{center}
\small
\begin{longtable}{@{}p{1.6cm}p{12.4cm}@{}}
\toprule
\textbf{ID} & \textbf{Functional requirement} \\
\midrule
\endhead
FR-01 & Institute-scoped user login with short-lived access sessions and
refresh rotation. \\
FR-02 & Multiple institutes served from one deployment with strict data
isolation per institute. \\
FR-03 & Institute admin creates users (teachers, students), assigns roles, and
manages role permissions from an explicit catalogue. \\
FR-04 & Academic structure: academic years, classes, divisions, subjects,
class--subject offerings, teacher assignments, and student placements. \\
FR-05 & Upload teaching material (PDF, image, document, or text) with chunked
upload for large files. \\
FR-06 & Distributed OCR processing: embedded-text-first extraction with
per-page PaddleOCR fallback, chunk leasing, and retry; per-page text
corrections. \\
FR-07 & AI syllabus analysis producing a confirmed global curriculum (chapters,
topics, units, objectives, learning outcomes). \\
FR-08 & AI generation of study content (notes, flashcards, Cornell notes,
summaries, important concepts) from eligible sources with provenance. \\
FR-09 & Question types as data (global starters plus institute-custom types). \\
FR-10 & Question bank CRUD with approval workflow and difficulty/source
tracking. \\
FR-11 & AI question-bank generation in batches with pattern-coverage checks
(COVERED, AWAITING\_APPROVAL, INSUFFICIENT, GENERATING). \\
FR-12 & Paper patterns: manual authoring, text/material analysis, deterministic
validation, and approval-to-lock. \\
FR-13 & Question papers: selection from pattern, coverage checking, generation
of missing items, and conversion to assessments. \\
FR-14 & Assessments: authoring, pattern-based selection, publishing,
activation, and completion with schedule validation. \\
FR-15 & Student attempts: timed start with an immutable question snapshot,
per-answer saving, and idempotent submission with server-side evaluation. \\
FR-16 & Aggregate analytics: score distribution, per-question accuracy, and
topic/difficulty performance. \\
FR-17 & Ungraded practice: flashcard sessions and question sessions with
instant per-item feedback and no score. \\
FR-18 & Export to PDF, DOCX, and XLSX for content, question sets, papers,
patterns, and results. \\
FR-19 & Asynchronous job lifecycle with retry, cancel, and status polling. \\
FR-20 & Health checks for the API and OCR service. \\
\bottomrule
\end{longtable}
\end{center}

\section{Non-functional Requirements}

\begin{center}
\small
\begin{longtable}{@{}p{1.6cm}p{12.4cm}@{}}
\toprule
\textbf{ID} & \textbf{Non-functional requirement} \\
\midrule
\endhead
NFR-01 & \textbf{Security.} Cookie-based sessions with httpOnly attributes,
double-submit CSRF on state-changing requests, origin check on login, and
throttling on authentication endpoints. \\
NFR-02 & \textbf{Tenant isolation.} Every query touching user data is scoped
by the active institute; cross-tenant access must fail by default of the
enforcement layers. \\
NFR-03 & \textbf{Authorization.} Permission keys are DB-fresh per request with
default-deny and \texttt{manage} implication; academic scope restricts
permission reach. \\
NFR-04 & \textbf{Availability of the API.} The API never blocks on long
operations; heavy work happens in workers and returns \texttt{202 Accepted}. \\
NFR-05 & \textbf{Scalability of OCR and AI.} OCR compute is distributed and
lease-based; AI generation has bounded concurrency to protect the model
gateway. \\
NFR-06 & \textbf{Maintainability.} Monorepo with TypeScript packages shared
between web and API; typed shared contracts; strict typing in both TypeScript
and Python codebases. \\
NFR-07 & \textbf{Portability.} Everything runs from one Docker Compose file
with no service exposed on a host port in production. \\
NFR-08 & \textbf{Traceability.} Generated AI artifacts store provenance
(source type, source id, revisions) and are marked AI-generated. \\
\bottomrule
\end{longtable}
\end{center}

\section{User Roles}

\noindent The platform defines three built-in membership roles and one
disjoint platform role:

\begin{itemize}
  \item \textbf{Institute admin (\texttt{INSTITUTE\_ADMIN}):} full management
  of an institute's structure and content; holds \texttt{manage} across the
  institute permission catalogue and whole-institute academic scope.
  \item \textbf{Teacher (\texttt{TEACHER}):} creates and manages teaching
  material, content, syllabus, questions, patterns, question papers, and
  assessments within their assigned subjects; reads results and analytics.
  \item \textbf{Student (\texttt{STUDENT}):} reads content and materials within
  their placement scope, takes assessments (creating attempts), and uses
  practice sessions.
  \item \textbf{Platform admin (\texttt{SUPER\_ADMIN}):} operates on the
  platform plane (institutes, OCR worker registry) through platform-domain
  permissions, disjoint from institute membership.
\end{itemize}

\section{Constraints and Assumptions}

\begin{itemize}
  \item \textbf{Data.} PostgreSQL is the only primary database; MongoDB and
  SQLite are not permitted.
  \item \textbf{AI access.} AI calls must go through an internal provider
  gateway; no direct cloud AI SDK is called from application code.
  \item \textbf{Delegation.} Only OCR and asynchronous AI generation are
  permitted to move off the request thread; the API remains a single
  modular monolith.
  \item \textbf{Provisioning.} Users are provisioned by the institute, not by
  public self-registration.
  \item \textbf{Text evaluation.} Only objective question formats are
  auto-graded; subjective answers are recorded but not graded.
  \item \textbf{Fonts and documents.} Exported assets must be deterministic and
  printable; PDF generation is browser-based.
\end{itemize}

\section{Environment of Use}

\noindent The system runs as a Docker Compose stack. Its runtime environment
comprises a PostgreSQL~17 database, RabbitMQ~3 message broker, Node.js~20
(TypeScript) API and web services, a Python~3.12 OCR service, Python workers, an
internal AI gateway, an nginx reverse proxy, and a Cloudflare Tunnel for the
only public ingress. No internal service is published on a host port in
production.